\chapter{Kohn--Sham theory on lattices}
\label{chap:ks}

Kohn--Sham theory replaces the interacting system by a non-interacting one with the same ground-state
density. On a lattice this requires (i) the structure of non-interacting ground states (aufbau), (ii) the
existence of a Kohn--Sham potential for a given density (non-interacting $v$-representability), and (iii) its
uniqueness. Statements without a check mark are open targets of the formalization.

\section{Non-interacting ground states}

\begin{definition}[Orbital energies]
  \label{def:orbEnergy}
  \lean{LatticeDFT.orbEnergy}
  \leanok
  For orbital energies $\varepsilon_a$ and an orbital subset $S$ of size $N$, $E_S=\sum_{a\in S}\varepsilon_a$.
\end{definition}

\begin{lemma}[Slater determinants of eigen-orbitals]
  \label{lem:slaterSub_eigen}
  \lean{LatticeDFT.slaterSub_eigen}
  \leanok
  \uses{thm:gamma_slater, def:orbEnergy}
  If $h\varphi_a=\varepsilon_a\varphi_a$, then $\Gamma(h)\Psi_S=E_S\Psi_S$.
\end{lemma}
\begin{proof}
  \leanok
  Theorem~\ref{thm:gamma_slater}.
\end{proof}

\begin{theorem}[Variational lower bound]
  \label{thm:qf_gamma_ge}
  \lean{LatticeDFT.slaterMat, LatticeDFT.slaterMat_orth, LatticeDFT.dot_slaterMat, LatticeDFT.qf_gamma_ge}
  \leanok
  \uses{thm:slater_orthonormal, lem:slaterSub_eigen}
  For an orthonormal eigenbasis of $h$ and every many-body state $\Psi$,
  $\ip{\Psi}{\Gamma(h)\Psi}\ge\big(\min_SE_S\big)\|\Psi\|^2$.
\end{theorem}
\begin{proof}
  \leanok
  The matrix $P$ whose columns are the $\Psi_S$ is orthogonal ($P^TP=1$, hence $PP^T=1$), and
  $\Gamma(h)P=P\,\diag(E_S)$. In the coefficients $c=P^T\Psi$ the energy is $\sum_SE_Sc_S^2$.
\end{proof}

\begin{theorem}[Aufbau]
  \label{thm:aufbau}
  \lean{LatticeDFT.exists_eigenbasis, LatticeDFT.conf_nonempty, LatticeDFT.aufbau}
  \leanok
  \uses{thm:qf_gamma_ge, def:IsGroundState}
  For a symmetric $h$ and $N\le M$, the ground-state energy of $\Gamma(h)$ is the smallest sum of $N$
  orbital energies. The Slater determinant of a minimising set of orbitals is a many-body ground state.
\end{theorem}
\begin{proof}
  \leanok
  Spectral theorem (Mathlib) for an orthonormal eigenbasis, then Theorem~\ref{thm:qf_gamma_ge}.
\end{proof}

\begin{theorem}[Lowest orbitals]
  \label{thm:aufbau_lowest}
  \lean{LatticeDFT.orbEnergy_eq_sum, LatticeDFT.sum_le_sum_of_forall_le, LatticeDFT.aufbau_lowest}
  \leanok
  \uses{thm:aufbau}
  A set $S$ minimises $E_S$ if and only if it consists of $N$ orbitals of lowest energy (with arbitrary
  choices among degenerate ones at the Fermi level).
\end{theorem}
\begin{proof}
  \leanok
  If an occupied $a$ lies above an unoccupied $b$, swapping them lowers $E_S$. Conversely, compare $S$ with
  any $T$. The orbitals $S\setminus T$ and $T\setminus S$ are equally many, and each element of the first lies
  below each element of the second, so $\sum_{S\setminus T}\varepsilon\le\sum_{T\setminus S}\varepsilon$.
\end{proof}

\section{Uniqueness of the Kohn--Sham potential}

\begin{theorem}[Generic uniqueness]
  \label{thm:ks_potential_generic_unique}
  \lean{LatticeDFT.ks_potential_generic_unique}
  \leanok
  \uses{thm:generic_hk_lattice_fermions}
  For spinless fermions on a connected graph with $t\ne0$, the Kohn--Sham potential reproducing a ground
  state is unique up to a constant outside a closed, nowhere dense, null set of potentials. It is not unique
  everywhere: the non-interacting two-fermion square has a whole plane of exceptions
  (Theorems~\ref{thm:hk_fails_square}, \ref{thm:square_wall}).
\end{theorem}
\begin{proof}
  \leanok
  Theorem~\ref{thm:generic_hk_lattice_fermions} with vanishing interaction.
\end{proof}

\begin{theorem}[Generic uniqueness, edge-dependent hopping]
  \label{thm:ks_potential_generic_unique_w}
  \lean{LatticeDFT.ks_potential_generic_unique_w}
  \leanok
  \uses{thm:generic_hk_lattice_fermions}
  The same holds for symmetric hopping amplitudes $t_{ab}$ on the edges, all of the same sign.
\end{theorem}
\begin{proof}
  \leanok
  The weighted generic Hohenberg--Kohn theorem with vanishing interaction.
\end{proof}

\section{Lieb's functional and convex duality}

Let $E(v)$ be the ground-state energy of $H(v)$. Ensemble densities of configurations form the density domain
$\mathcal{D}=\{\sum_c\lambda_cn(c):\lambda\ \text{a probability vector}\}$.

\begin{definition}[Lieb's functional]
  \label{def:lieb}
  \lean{LatticeDFT.Egs, LatticeDFT.DensDomain, LatticeDFT.lieb}
  \leanok
  \uses{def:hamN, thm:exists_groundState}
  $F(\rho)=\sup_v\,[E(v)-v\cdot\rho]$, the Legendre transform of $E$.
\end{definition}

\begin{lemma}
  \label{lem:lieb_basic}
  \lean{LatticeDFT.Egs_le_config, LatticeDFT.lieb_bdd, LatticeDFT.le_lieb, LatticeDFT.dens_mem_domain}
  \leanok
  \uses{def:lieb}
  $E(v)\le (H_0)_{cc}+\sum_iv_in_i(c)$ for every configuration $c$. Hence $F$ is finite on $\mathcal{D}$, and
  ground-state densities lie in $\mathcal{D}$.
\end{lemma}
\begin{proof}
  \leanok
  Use the configuration $e_c$ as a trial state, and take convex combinations.
\end{proof}

\begin{theorem}[Lieb duality]
  \label{thm:lieb_duality}
  \lean{LatticeDFT.lieb_duality}
  \leanok
  \uses{lem:lieb_basic, thm:exists_groundState}
  $E(v)=\min_{\rho\in\mathcal{D}}\,[F(\rho)+v\cdot\rho]$, the minimum being attained at the density of any
  ground state of $H(v)$.
\end{theorem}
\begin{proof}
  \leanok
  ``$\le$'' holds by definition of $F$. For a ground state $\psi$, the variational principle gives
  $E(v')-v'\cdot\rho_\psi\le\qf(H_0,\psi)=E(v)-v\cdot\rho_\psi$ for all $v'$.
\end{proof}

\begin{theorem}[Convexity]
  \label{thm:lieb_convexOn}
  \lean{LatticeDFT.lieb_convexOn}
  \leanok
  \uses{lem:lieb_basic}
  $F$ is convex on the (convex) density domain.
\end{theorem}
\begin{proof}
  \leanok
  A supremum of affine functions.
\end{proof}

\section{Existence of the Kohn--Sham potential}

Assume nonnegative occupations with fixed particle number ($\sum_in_i(c)=N$), rigid occupations, and a
density $\rho=\sum_c\lambda_cn(c)$ with all $\lambda_c>0$ (a strictly positive mixture).

\begin{lemma}[Regularity of $E$]
  \label{lem:Egs_regular}
  \lean{LatticeDFT.sum_dens, LatticeDFT.qf_hamN_shift, LatticeDFT.Egs_shift, LatticeDFT.Nr_nonneg, LatticeDFT.qf_hamN_diff_le, LatticeDFT.Egs_lipschitz, LatticeDFT.continuous_Egs}
  \leanok
  \uses{lem:lieb_basic}
  $E(v+k)=E(v)+kN$, and $|E(v)-E(v')|\le N\|v-v'\|_\infty$. In particular $E$ is continuous.
\end{lemma}
\begin{proof}
  \leanok
  Compare expectation values in the ground states of both Hamiltonians.
\end{proof}

\begin{theorem}[Coercivity and a maximiser]
  \label{thm:exists_maximiser}
  \lean{LatticeDFT.aval, LatticeDFT.mval, LatticeDFT.mval_le, LatticeDFT.continuous_mval, LatticeDFT.mval_smul, LatticeDFT.weighted_aval, LatticeDFT.mval_nonpos, LatticeDFT.const_of_mval_eq_zero, LatticeDFT.coercive, LatticeDFT.sum_rho, LatticeDFT.exists_maximiser}
  \leanok
  \uses{lem:Egs_regular, def:NonUV}
  There is $\delta>0$ with $E(v)-v\cdot\rho\le K-\delta\|v\|$ whenever $\sum_iv_i=0$. Consequently
  $v\mapsto E(v)-v\cdot\rho$ attains its maximum.
\end{theorem}
\begin{proof}
  \leanok
  Let $m(v)=\min_c\,(v\cdot n(c)-v\cdot\rho)$. Since $\sum_c\lambda_c(v\cdot n(c)-v\cdot\rho)=0$ with
  $\lambda>0$, we have $m\le0$, with equality only for constant $v$ (rigidity). Compactness of the unit sphere
  in $\{\sum_iv_i=0\}$ gives the uniform $\delta$. Then maximise over a large ball and use shift invariance.
\end{proof}

\begin{theorem}[Ground states below the target]
  \label{thm:exists_gs_below}
  \lean{LatticeDFT.qf_hamN_add, LatticeDFT.continuous_qf, LatticeDFT.continuous_dens, LatticeDFT.isCompact_sphere1, LatticeDFT.exists_gs_below}
  \leanok
  \uses{thm:exists_maximiser, thm:exists_groundState}
  At a maximiser $v^*$, for every direction $w$ some normalised ground state $\psi$ of $H(v^*)$ satisfies
  $w\cdot\rho_\psi\le w\cdot\rho$.
\end{theorem}
\begin{proof}
  \leanok
  Ground states $\psi_k$ of $H(v^*+t_kw)$ satisfy $w\cdot\rho_{\psi_k}\le w\cdot\rho$ by maximality.
  A convergent subsequence (the unit sphere is compact) gives a ground state of $H(v^*)$, by continuity of $E$.
\end{proof}

\begin{definition}[Ground-state ensembles]
  \label{def:Ens}
  \lean{LatticeDFT.Ens, LatticeDFT.ensDens, LatticeDFT.ensDens_vecMulVec, LatticeDFT.pure_mem_Ens}
  \leanok
  \uses{thm:exists_groundState}
  $\mathrm{Ens}(v)$ is the set of density matrices $\Gamma$ (symmetric, positive semidefinite, trace one)
  with $H(v)\Gamma=E(v)\Gamma$, i.e.\ supported on the ground-state space. Their density is
  $\rho_\Gamma(i)=\sum_cn_i(c)\Gamma_{cc}$. Pure ground states $\psi\psi^T$ belong to $\mathrm{Ens}(v)$.
\end{definition}

\begin{lemma}[Ensembles form a compact convex set]
  \label{lem:Ens_compact}
  \lean{LatticeDFT.quad_single, LatticeDFT.Ens_bound, LatticeDFT.convex_Ens, LatticeDFT.isCompact_Ens, LatticeDFT.continuous_ensDens}
  \leanok
  \uses{def:Ens}
  $\mathrm{Ens}(v)$ is convex and compact (closed, with entries in $[-1,1]$).
\end{lemma}
\begin{proof}
  \leanok
  Entry bounds follow from positivity in the directions $e_a\pm e_b$ and the trace condition.
\end{proof}

\begin{theorem}[Ensemble $v$-representability]
  \label{thm:ensemble_vrep}
  \lean{LatticeDFT.ensemble_v_representable}
  \leanok
  \uses{thm:exists_gs_below, lem:Ens_compact}
  Every strictly positive mixture $\rho$ of configuration densities is the density of a ground-state ensemble
  of $H(v)$ for some $v$. This holds for any symmetric $H_0$, interacting or not.
\end{theorem}
\begin{proof}
  \leanok
  Take a maximiser $v^*$. If $\rho$ were not an ensemble density, the compact convex set of ensemble
  densities could be separated from $\rho$ by a direction $w$ (Hahn--Banach). That contradicts
  Theorem~\ref{thm:exists_gs_below}.
\end{proof}

\begin{theorem}[Kohn--Sham potentials exist]
  \label{thm:ks_potential_exists}
  \lean{LatticeDFT.sum_occF, LatticeDFT.v_representable_lattice_fermions, LatticeDFT.ks_potential_exists}
  \leanok
  \uses{thm:ensemble_vrep, thm:rigid_occF, def:H0F}
  For $N$ spinless fermions ($0<N<M$) on any graph and any density--density interaction, every strictly
  positive mixture of configuration densities is ensemble $v$-representable. With vanishing interaction this
  is the existence of a Kohn--Sham potential.
\end{theorem}
\begin{proof}
  \leanok
  Theorem~\ref{thm:ensemble_vrep} with occupations $n_i(s)=[i\in s]$.
\end{proof}

\begin{lemma}[Vertex description of the hypersimplex]
  \label{lem:hypersimplex_mixture}
  \lean{LatticeDFT.IsMixture, LatticeDFT.nfrac, LatticeDFT.mixture_combo, LatticeDFT.hypersimplex_mixture}
  \leanok
  Every $\rho\in[0,1]^M$ with $\sum_i\rho_i=N$ is a mixture (convex combination) of $N$-particle
  configuration densities.
\end{lemma}
\begin{proof}
  \leanok
  Induction on the number of fractional sites ($0<\rho_i<1$). If there are none, $\rho$ is the occupation
  vector of a configuration. Since $\sum_i\rho_i$ is an integer, there is never exactly one fractional site,
  so otherwise there are two, $i\neq j$. Moving particle weight from $j$ to $i$ (or back) as far as possible
  gives two densities $\rho^\pm$, still in the box with the same sum, each with fewer fractional sites, and
  $\rho$ is a convex combination of them. A convex combination of mixtures is a mixture.
\end{proof}

\begin{theorem}[Interior densities]
  \label{thm:interior_positive_mixture}
  \lean{LatticeDFT.interior_positive_mixture}
  \leanok
  \uses{lem:hypersimplex_mixture}
  A density with $0<\rho_i<1$ and $\sum_i\rho_i=N$ is a strictly positive mixture of the $N$-particle
  configuration densities.
\end{theorem}
\begin{proof}
  \leanok
  Let $\bar\rho$ be the uniform mixture of all configurations and $\varepsilon=\min_i\min(\rho_i,1-\rho_i)>0$.
  Then $\rho'=(\rho-\varepsilon\bar\rho)/(1-\varepsilon)$ lies in the hypersimplex, so it is a mixture
  $\sum_s\mu_s n_s$ (Lemma~\ref{lem:hypersimplex_mixture}). The weights $(1-\varepsilon)\mu_s+\varepsilon/|\Conf|$
  are strictly positive and reproduce $\rho$.
\end{proof}

\begin{theorem}[Existence of Kohn--Sham potentials, textbook form]
  \label{thm:ks_potential_exists_interior}
  \lean{LatticeDFT.v_representable_interior, LatticeDFT.ks_potential_exists_interior}
  \leanok
  \uses{thm:interior_positive_mixture, thm:ks_potential_exists}
  For $N$ spinless fermions with $0<N<M$ on any graph, any hopping $t$ and any configuration-diagonal
  interaction, every density with $0<\rho_i<1$ and $\sum_i\rho_i=N$ is the density of a ground-state ensemble
  of $H(v)$ for some potential $v$. In particular (no interaction), every such density has a Kohn--Sham
  potential.
\end{theorem}
\begin{proof}
  \leanok
  Combine Theorem~\ref{thm:interior_positive_mixture} with Theorem~\ref{thm:ks_potential_exists}.
\end{proof}

\section{The regularised Kohn--Sham iteration}

We follow Penz, Laestadius, Tellgren and Ruggenthaler, Phys.~Rev.~Lett.~123, 037401 (2019), together with
its erratum (with P.~E.~Lammert). The original letter claimed that the iterates converge. The erratum
points out that $\|x_{i+1}-x_i\|\to0$ does not imply this, and proves the weaker statement formalised here:
the iterates are bounded, and every accumulation point is a ground-state quasi-density.

Densities and potentials live in $X=\ell^2(V)$ with inner product $\ip{\cdot}{\cdot}$; $\varepsilon>0$ is
the regularisation parameter.

\subsection{An abstract convergence theorem}

\begin{definition}[Moreau--Yosida smoothness]
  \label{def:IsMYSmooth}
  \lean{LatticeDFT.IsMYSmooth}
  \leanok
  $g:X\to X$ is an $\varepsilon$-gradient of $F:X\to\R$ if, for all $x,y$,
  \[F(x)+\ip{g(x)}{y-x}\le F(y)\le F(x)+\ip{g(x)}{y-x}+\tfrac1{2\varepsilon}\|y-x\|^2 .\]
\end{definition}

\begin{lemma}[Co-coercivity]
  \label{lem:cocoercive}
  \lean{LatticeDFT.IsMYSmooth.add_linear, LatticeDFT.IsMYSmooth.lower_strong, LatticeDFT.IsMYSmooth.cocoercive, LatticeDFT.IsMYSmooth.norm_sub_le, LatticeDFT.IsMYSmooth.continuous, LatticeDFT.IsMYSmooth.norm_le, LatticeDFT.IsMYSmooth.continuous_fun}
  \leanok
  \uses{def:IsMYSmooth}
  The structure is preserved by adding a linear term. Moreover
  $\varepsilon\|g(x)-g(y)\|^2\le\ip{g(x)-g(y)}{x-y}$, so $g$ is $\varepsilon^{-1}$-Lipschitz, and $F$ is
  continuous.
\end{lemma}
\begin{proof}
  \leanok
  Compare $x$, $y$ and $z=y-\varepsilon(g(y)-g(x))$ using both inequalities, then symmetrise.
\end{proof}

\begin{lemma}[Regularised ground states exist]
  \label{lem:exists_stationary}
  \lean{LatticeDFT.IsMYSmooth.exists_stationary}
  \leanok
  \uses{lem:cocoercive}
  If the sublevel sets of $F+\ip{w}{\cdot}$ are bounded, then there is a $y$ with $g(y)+w=0$.
\end{lemma}
\begin{proof}
  \leanok
  A minimiser exists by compactness. At a minimiser the upper bound in direction $-\varepsilon(g(y)+w)$
  forces $g(y)+w=0$.
\end{proof}

\begin{definition}[MYKSODA]
  \label{def:myksoda}
  \lean{LatticeDFT.ksStep, LatticeDFT.ksStep_mul}
  \leanok
  Let $F$ (interacting) and $T$ (non-interacting reference) have $\varepsilon$-gradients $g_F$, $g_T$, and let
  $v$ be the external potential. The iteration is:
  (a) $w_i = v+g_F(x_i)-g_T(x_i)$ (Kohn--Sham potential);
  (b) choose $x_i'$ with $g_T(x_i')=-w_i$ (ground state of the regularised reference system);
  (c) $x_{i+1}=x_i+s_i(x_i'-x_i)$ with $s_i=-\varepsilon\ip{g_F(x_i)+v}{d_i}/\|d_i\|^2$, where
  $d_i = x_i'-x_i$ (and $s_i=0$ if $d_i=0$).
\end{definition}

\begin{lemma}[One step]
  \label{lem:step_descent}
  \lean{LatticeDFT.step_descent, LatticeDFT.step_lower}
  \leanok
  \uses{def:myksoda, def:IsMYSmooth}
  With $G=F+\ip{v}{\cdot}$ and $h_i=x_{i+1}-x_i$: $G(x_{i+1})\le G(x_i)-\|h_i\|^2/(2\varepsilon)$. If
  $\ip{g}{d}\le-\varepsilon\|g\|^2$, then $\varepsilon^2\|g\|^2\le\|h\|\,\|d\|$.
\end{lemma}
\begin{proof}
  \leanok
  Insert the step into the upper bound; the chosen step is the vertex of the parabola.
\end{proof}

\begin{theorem}[Convergence of MYKSODA, abstract form]
  \label{thm:myksoda_converges}
  \lean{LatticeDFT.myksoda_converges, LatticeDFT.myksoda_exists}
  \leanok
  \uses{lem:step_descent, lem:cocoercive, lem:exists_stationary}
  Assume that the proximal points $y-\varepsilon g_T(y)$ are bounded and the sublevel sets of $G$ are
  bounded. Then the iteration exists from every starting point, and:
  \begin{enumerate}
    \item the energies $G(x_i)$ decrease;
    \item $g_F(x_i)+v\to0$;
    \item there is an accumulation point;
    \item every accumulation point $z=\lim x_{\varphi(i)}$ satisfies $g_F(z)+v=0$ and minimises $G$;
      $G(x_i)\to G(z)$; and $w_{\varphi(i)}\to-g_T(z)$.
  \end{enumerate}
\end{theorem}
\begin{proof}
  \leanok
  Co-coercivity of $g_T$ applied to $x_i$ and $x_i'$ gives the key estimate
  $\ip{g_F(x_i)+v}{d_i}\le-\varepsilon\|g_F(x_i)+v\|^2$ (eq.~(17) of the paper).
  \begin{itemize}
    \item The energies decrease and are bounded below on the bounded sublevel set, so
      $\sum_i\|h_i\|^2<\infty$.
    \item The directions $d_i$ are bounded, so $\|g_F(x_i)+v\|^2\le\|h_i\|\,D/\varepsilon^2\to0$.
    \item Bolzano--Weierstrass gives an accumulation point; continuity of the gradients gives the limits.
    \item Convexity gives $G(x_i)\le G(z)+\|g_F(x_i)+v\|\,\|z-x_i\|$, hence the energy limit.
  \end{itemize}
\end{proof}

\subsection{Moreau--Yosida envelopes}

\begin{definition}[Moreau envelope]
  \label{def:myEnv}
  \lean{LatticeDFT.prox, LatticeDFT.myEnv, LatticeDFT.myGrad}
  \leanok
  For $f$ convex and lower semicontinuous on a nonempty compact convex $D\subset X$ (and $+\infty$ outside):
  \begin{itemize}
    \item $\mathrm{prox}(x)$ minimises $f(q)+\|x-q\|^2/(2\varepsilon)$ over $q\in D$;
    \item $f_\varepsilon(x)$ is the minimum value;
    \item $\nabla f_\varepsilon(x)=\varepsilon^{-1}(x-\mathrm{prox}(x))$.
  \end{itemize}
\end{definition}

\begin{lemma}[Proximal point]
  \label{lem:prox}
  \lean{LatticeDFT.prox_spec, LatticeDFT.nonneg_of_forall_small, LatticeDFT.prox_vi}
  \leanok
  \uses{def:myEnv}
  The proximal point exists and lies in $D$. It satisfies the variational inequality
  $f(p)+\varepsilon^{-1}\ip{x-p}{q-p}\le f(q)$ for all $q\in D$.
\end{lemma}
\begin{proof}
  \leanok
  A lower semicontinuous function attains its minimum on a compact set. For the inequality, compare $p$ with
  $p+t(q-p)$ and let $t\to0^+$.
\end{proof}

\begin{theorem}[Moreau envelopes are MY-smooth]
  \label{thm:isMYSmooth_myEnv}
  \lean{LatticeDFT.isMYSmooth_myEnv, LatticeDFT.sub_smul_myGrad, LatticeDFT.myEnv_prox_bounded, LatticeDFT.myEnv_sublevel_bounded}
  \leanok
  \uses{lem:prox, def:IsMYSmooth}
  $\nabla f_\varepsilon$ is an $\varepsilon$-gradient of $f_\varepsilon$. The proximal points are bounded,
  and the sublevel sets of $f_\varepsilon+\ip{w}{\cdot}$ are bounded for every $w$.
\end{theorem}
\begin{proof}
  \leanok
  The upper bound uses $\mathrm{prox}(x)$ as a competitor at $y$. The lower bound follows from the
  variational inequality plus $\|u-b\|^2\ge0$. Growth: $f_\varepsilon(y)\ge\min f+\mathrm{dist}(y,D)^2/(2\varepsilon)$.
\end{proof}

\subsection{The lattice iteration}

\begin{lemma}[Lieb's functional on $\ell^2$]
  \label{lem:liebE}
  \lean{LatticeDFT.densDomE, LatticeDFT.liebE, LatticeDFT.Freg, LatticeDFT.gFreg, LatticeDFT.inner_eq_dot, LatticeDFT.densDomain_eq, LatticeDFT.isCompact_densDomain, LatticeDFT.isCompact_densDomE, LatticeDFT.densDomE_nonempty, LatticeDFT.convexOn_liebE, LatticeDFT.convex_densDomE, LatticeDFT.liebE_lsc, LatticeDFT.isMYSmooth_Freg}
  \leanok
  \uses{thm:lieb_duality, thm:isMYSmooth_myEnv}
  The density domain is a nonempty compact convex set: the image of the standard simplex. Lieb's functional
  is convex on it. It is also lower semicontinuous there, being a supremum of continuous functions. Hence
  its Moreau--Yosida regularisation $F_\varepsilon$ is MY-smooth.
\end{lemma}
\begin{proof}
  \leanok
  Combine the preceding results.
\end{proof}

\begin{theorem}[Back to densities]
  \label{thm:density_of_quasi}
  \lean{LatticeDFT.density_of_quasi}
  \leanok
  \uses{lem:liebE, lem:prox}
  If $\nabla F_\varepsilon(z)=-v$, then $\rho=z+\varepsilon v$ is an ensemble density and
  $F(\rho)+v\cdot\rho=E(v)$: $\rho$ is a ground-state density in the sense of Lieb's variational
  principle.
\end{theorem}
\begin{proof}
  \leanok
  $\rho=\mathrm{prox}(z)$. The variational inequality shows that $\rho$ minimises $F+v\cdot$ over the domain,
  and Lieb duality identifies the minimum with $E(v)$.
\end{proof}

\begin{theorem}[Convergence of the regularised Kohn--Sham iteration on a lattice]
  \label{thm:ks_iteration}
  \lean{LatticeDFT.lattice_myksoda, LatticeDFT.lattice_myksoda_exists}
  \leanok
  \uses{thm:myksoda_converges, thm:density_of_quasi, lem:liebE}
  Let $H_0$ (interacting) and $H_0'$ (reference) be symmetric, $v\in\ell^2(V)$ and $\varepsilon>0$. The
  MYKSODA iteration with the regularised Lieb functionals of both systems exists from every starting point.
  Its energies decrease, and it has accumulation points. For every accumulation point $z$:
  \begin{itemize}
    \item $z$ minimises the regularised energy, and the energies converge to that minimum;
    \item $\rho=z+\varepsilon v$ is a ground-state density of the interacting system;
    \item the Kohn--Sham potentials converge along the subsequence to $v_{KS}$, and
      $\rho_{KS}=z+\varepsilon v_{KS}$ is a ground-state density of the reference system in $v_{KS}$
      (eq.~(24) of the paper).
  \end{itemize}
\end{theorem}
\begin{proof}
  \leanok
  Instantiate Theorem~\ref{thm:myksoda_converges} with $F_\varepsilon$ and $T_\varepsilon$, then apply
  Theorem~\ref{thm:density_of_quasi} to both systems.
\end{proof}
